\clearpage\fancyhead[L]{\small PALLAS.}\hypertarget{native-336-lower}{}\label{native-336-lower}
\begin{center}\textbf{PALLAS.}\par\vspace{4pt}\hrule\vspace{4pt}\fontsize{10}{12}\selectfont Astronomisches Jahrbuch for 1814. pp. 246\,..\,249. Berlin 1811.\par\vspace{4pt}\hrule\vspace{4pt}\end{center}

\begin{flushright}Göttingen, 10 February 1811.\end{flushright}
\par\indent On 14 January of this year I had a short report on the perturbations of Pallas by Jupiter printed in our learned Notices. Of course this is still short of the \emph{general} theory of those perturbations, which will take much more work if it is to do everything one could wish for. I have already made a start on it, though there is no hurry yet, because I can carry on treating the perturbations as I did in my recent work for many more years. As soon as this year’s opposition has been observed, I will carry on with the work, perhaps then treating Juno and Vesta in a similar way as well. I ask you to observe the coming $\text{\AstroSymbols ☍}$ of Pallas carefully. To make this easier, here is the improved ephemeris, since the agreement of my new results has been fully confirmed.
\Needspace{14\baselineskip}
\begin{center}\fontsize{10}{12}\selectfont\setlength{\tabcolsep}{4pt}\renewcommand{\arraystretch}{1.15}\begin{tabular}{lll}
Feb. 11 & R.A. $148^\circ\,3'$ & Decl. $15^\circ\,43'$ S.\\
\phantom{Feb. }15 & $147\;\;20$ & $14\;\;16$\\
\phantom{Feb. }19 & $146\;\;37$ & $12\;\;42$\\
\phantom{Feb. }23 & $145\;\;55$ & $11\;\;3$\\
\phantom{Feb. }27 & $145\;\;17$ & $9\;\;19$\\
March \phantom{1}3 & R.A. $144^\circ\,43'$ & Decl. $7^\circ\,32'$ S.\\
\phantom{March }7 & $144\;\;13$ & $5\;\;44$\\
11 & $143\;\;49$ & $3\;\;57$\\
15 & $143\;\;31$ & $2\;\;11$\\
19 & $143\;\;20$ & $0\;\;29$\\
\end{tabular}\end{center}

\SourceLeaf{337}\Needspace{12\baselineskip}\par\indent Professor \textsc{Harding} made the following observations of Pallas in Mannheim.
\begin{center}\fontsize{10}{12}\selectfont\setlength{\tabcolsep}{4pt}\renewcommand{\arraystretch}{1.15}
\begin{tabular}{l|c|c|c}
1811 January 20 & $14^{\mathrm u}\,6^{\mathrm m}\,38^{\mathrm s}6$ & $151^\circ\,10'\,36''2$ & $-21^\circ\,7'\,39''7$ \\
\phantom{1811 January }21 & $14\;\;2\;\;18.1$ & $151\;\;4\;\;53.5$ & $-20\;\;57\;\;27.7$ \\
\phantom{1811 January }22 & $13\;\;57\;\;58.4$ & & $-20\;\;47\;\;51.7$ \\
\phantom{1811 January }23 & $13\;\;53\;\;38.6$ & $150\;\;52\;\;48.2$ & $-20\;\;37\;\;58.8$ \\
\phantom{1811 January }24 & $13\;\;49\;\;16.4$ & $150\;\;46\;\;0.6$ & $-20\;\;27\;\;9.7$ \\
\phantom{1811 January }25 & $13\;\;44\;\;53.0$ & $150\;\;39\;\;9.9$ & $-20\;\;16\;\;9.4$ \\
\end{tabular}
\end{center}


\par\indent The planet is now fairly bright, though Ceres, which is very close to the position given in its ephemeris, seemed to me to be a little brighter still.
\Needspace{18\baselineskip}\begin{flushright}Göttingen, 29 August 1811.\end{flushright}
\par\indent My dear friend, I have the pleasure of sending you the enclosed ephemerides for the next apparition of \emph{Pallas}, \emph{Juno} and \emph{Vesta}. Three skilled pupils of mine shared the calculation. You will find more about Pallas in issues 8 and 130 of the Göttingische Gelehrte Anzeigen [p. 327 of this volume]. Under my supervision Mr. \textsc{Nicolai} has now carried the calculation of Jupiter’s perturbation of Pallas forward another year, and he is now calculating the ephemeris for 1813. Further details about Juno are in issue 92 of our learned Notices [p. 325 of this volume.] The elements of Vesta are still provisional and will need a small improvement. In the meantime, here are the observations I made here with the mural quadrant.
\Needspace{9\baselineskip}\begin{center}\fontsize{10}{12}\selectfont\setlength{\tabcolsep}{4pt}\renewcommand{\arraystretch}{1.15}
\begin{tabular}{l|c|c|c}
1811 May 20 & $12^{\mathrm u}\,27^{\mathrm m}\,45''1$ M.\,T. & $244^\circ\,42'\,0''9$ & $-12^\circ\,28'\,10''1$ \\
\phantom{1811 May }24 & $12\;\;8\;\;3.4$ & $243\;\;42\;\;26.1$ & $-12\;\;30\;\;1.8$ \\
\phantom{1811 May }25 & $12\;\;3\;\;7.6$ & $243\;\;27\;\;25.3$ & $-12\;\;30\;\;47.8$ \\
\phantom{1811 May }29 & $11\;\;43\;\;22.9$ & $242\;\;27\;\;0.7$ & $-12\;\;35$ \\
\phantom{1811 May }30 & $11\;\;38\;\;27.3$ & $242\;\;12\;\;3.0$ & $-12\;\;35\;\;52.8$ \\
\end{tabular}
\end{center}


\par\indent From these observations, all of which I consider good, I derived the opposition as follows. 1811 May 15. $12^{\mathrm u}\,44^{\mathrm m}\,30^{\mathrm s}$ mean time in Göttingen, $243^\circ\,48'\,43''9$ true longitude, $8^\circ\,33'\,59''5$ geocentric latitude, north. As for the changes that Dr. \textsc{Triesnecker} thought he should make to my results for some oppositions of Ceres, as reported in the Yearbook for 1813, I would make the following remarks.

\par\indent 1) As far as I have checked the calculation of the opposition for 1807 from the Milan observations, I have found no errors and must stand by my result, 3 May $8^{\mathrm u}\,42^{\mathrm m}\,28^{\mathrm s}$ Paris time and $222^\circ\,14'\,9''1$ longitude. I will nevertheless repeat the whole calculation as soon as I have a chance to calculate the elements from the two oppositions of 1809 and 1811 (which have yet to be used).

\par\indent 2) Dr. \textsc{Triesnecker}’s conjecture that the geocentric latitude at the 1808 opposition should be changed from $12^\circ\,43'\,44''8$ to $12^\circ\,45'\,44''8$ is quite inadmissible. The longitude of the line of nodes likewise has to remain $80^\circ\,53'\,23''$ and must on no account be taken to be $5'$ greater. The sole reason Mr. Tr. was led to this conjecture was that with my perturbation tables the elements would otherwise agree less well with the latitude. But I am very surprised that he overlooked what I expressly pointed out. In calculating the latitudes I used the complete perturbation equations I later developed in place of my old tables. They have yet to appear in print, so I give them to you here as they were calculated in 1805.
\SourceLeaf{338}\Needspace{24\baselineskip}\begin{center}\emph{Perturbations of the latitude of Ceres by Jupiter.}

($\JupiterSymbol$ mean longitude of $\JupiterSymbol$, $\CeresSymbol$ mean longitude of $\CeresSymbol$)
\end{center}
\[
\begin{array}{l}
-9''65\,\sin(\JupiterSymbol-74^\circ\,13')\\
-14.17\,\sin(\CeresSymbol-2\JupiterSymbol+78^\circ\,35')\\
+27.47\,\sin(2\CeresSymbol-3\JupiterSymbol+78\;35)\\
-4.62\,\sin(3\CeresSymbol-4\JupiterSymbol+78\;35)\\
+5.51\,\sin(2\CeresSymbol-\JupiterSymbol-78\;35)\\
+0.99\,\sin(3\CeresSymbol-2\JupiterSymbol-78\;35)\\
+1.91
\end{array}
\]
\[
\begin{array}{l}
+8.61\,\sin(\CeresSymbol-\JupiterSymbol+95\;26)\\
-10.31\,\sin(2\CeresSymbol-2\JupiterSymbol+94\;10)\\
-6''16\,\sin(3\CeresSymbol-3\JupiterSymbol+112^\circ\,26')\\
-13.69\,\sin(2\JupiterSymbol-57^\circ\,21')\\
-6.77\,\sin(3\JupiterSymbol-\CeresSymbol-61\;25)\\
+4.30\,\sin(2\CeresSymbol-4\JupiterSymbol+72\;24)\\
-22.51\,\sin(3\CeresSymbol-5\JupiterSymbol+70\;27)\\
+1.65\,\sin(3\JupiterSymbol-56^\circ\,9')
\end{array}
\]
\par\indent The first six equations are the old ones with improvements, whereas the remaining nine, which can amount to one minute, are altogether new additions. It is therefore hardly surprising that taking them all into account changed the position of the node by $5'$. — These 15 latitude equations are still bearable, but for Pallas I have over 40 that exceed $1''$.\par
\indent We first saw the comet again on 22 August, but Professor \textsc{Harding} only measured some of its distances from fixed stars this morning. They have yet to be reduced.
